% Kapitel 1: Einleitung
\section{Einleitung}
\label{sec:einleitung}

\subsection{Motivation}
\label{sec:motivation}

Die Erlebniswelt von Zotter Schokolade zählt mit ihren über 300.000 jährlichen Besuchern (Stand 2025) zu den Top-3 der steirischen Sehenswürdigkeiten mit Bezahl-Eintritt \cite{steiermark_ausflugsziele}.
Diese enormen Besucherströme stellen eine nicht unbedeutende organisatorische Herausforderung für Mitarbeiter und Management dar.
Trotz der hohen Frequenz an Besuchern glich der Parkplatz aus datentechnischer Sicht einer "`Black Box"', da keine automatisierte Erfassung oder Auswertung von Fahrzeugbewegungen existierte.
Informationen über die tatsächliche Auslastung, Stoßzeiten oder die Verweildauer basierten auf subjektiven Beobachtungen statt auf Daten.
Dies erschwerte sowohl die operative Ressourcenplanung als auch strategische Entscheidungen erheblich. \\
Diese Diplomarbeit zielt darauf ab, diese Lücke zu schließen, indem ein modernes Kennzeichenerkennungssystem (ALPR) eingeführt wird, um Kennzeichen und Fahrzeugdaten zu erfassen, diese unter anderem auf Herkunft zu analysieren und anonymisiert zu speichern.
Auf diese Daten aufbauend, können Analysen, grafische Dashboards mit den wichtigsten Informationen und langfristige Vorhersagen getroffen werden.
Ziel ist es, mit diesen Möglichkeiten der Datensammlung und Analyse die strategische Planung, das Marketing und die Parkplatz-Logistik zu unterstützen.

\newpage
\subsection{Zielsetzung}
\label{sec:zielsetzung}

Das Ziel dieser Diplomarbeit ist es, mit Daten des Besucherparkplatzes der Zotter Schokolade GmbH Analysen und Visualisierungen zu bieten, mit deren Hilfe bessere Entscheidungen in den Bereichen Betriebsführung, Marketing und Ressourcenplanung ermöglicht werden.

Das entwickelte System soll:
\begin{itemize}
    \item Transparenz durch Echtzeit-Einblicke in Parkplatzauslastung, Besucherfrequenz und Stoßzeiten liefern
    \item Die Ressourcenplanung durch eine Datengrundlage für die strategische Planung unterstützen
    \item Das zielgerichtete Marketing durch die Bereitstellung von Herkunftsanalysen verbessern
    \item Eine Basis für die GHG (Greenhouse Gas Protocol) Scope 3 Emissionsberechnung des Besucheraufkommens liefern \cite{ghg_scope3}
    \item Datenschutz und DSGVO-Konformität durch Anonymisierung aller personenbezogenen Daten gewährleisten
\end{itemize}

Das System soll nicht:
\begin{itemize}
    \item Individuelle Besucher identifizieren oder tracken
    \item Bewegungsprofile einzelner Fahrzeuge über den Parkplatzkontext hinaus erstellen
    \item Personenbezogene Daten dauerhaft speichern
\end{itemize}

Durch die Abgrenzung dieser Ziele wird sichergestellt, dass das System nur als Werkzeug für betriebliche Optimierung dient, ohne dabei die Privatsphäre der Besucher zu gefährden.

\newpage
\subsection{Aufgabenstellung}
\label{sec:aufgabenstellung}

Um diesen Anforderungen gerecht zu werden, wurden sie in folgende Teilaufgaben aufgeteilt:
\begin{itemize}
    \item Installation bzw. Montage und Konfiguration der erforderlichen Hardware
    \item Implementierung eines ALPR-Systems (Automatic License Plate Recognition) zur Kennzeichenerkennung
    \item Erkennung von Ein- und Ausfahrten mit Zeitstempeln
    \item Auswertung der Kennzeichen nach Herkunftsland und Region mit der Erkennung von Fahrzeugtypen (PKW, LKW, etc.)
    \item DSGVO-konforme Anonymisierung und Speicherung der erfassten Daten
    \item Erstellung eines Scripts für die automatische Datenbank-Initialisierung
    \item Entwicklung einer automatisierten Backup- und Wiederherstellungslösung
    \item Implementierung eines Dashboards zur Visualisierung der Daten
    \item Integration eines E-Mail-Benachrichtigungssystems
    \item Bereitstellung einer Entwickler-Dokumentation für zukünftige Wartung und Erweiterungen
\end{itemize}

\newpage
\subsection{Aufbau der Arbeit}
\label{sec:aufbau}

Diese Diplomarbeit gliedert sich in Hauptkapitel, die den Entwicklungsprozess des Systems von den theoretischen Grundlagen bis zur Implementierung nachvollziehbar machen sollen.

Das nächste Kapitel befasst sich mit den theoretischen Grundlagen. Die eingesetzten Technologien (Microservices-Architektur, Containerisierung und ALPR-Technologie) werden näher beleuchtet und deren Relevanz für das Projekt begründet.

In Kapitel~\ref{sec:hardware} werden der Prozess der Hardwareauswahl, die eingesetzte Systemumgebung inklusive der getroffenen Entscheidungen bezüglich des Technologie-Stacks sowie die lokalen Standortgegebenheiten beschrieben.
Es werden die spezifischen Anforderungen an Kamera und Server unter Berücksichtigung von Umgebungsbedingungen und Vorgaben der Betreuerfirma evaluiert und die getroffenen Entscheidungen argumentiert.

Kapitel~\ref{sec:systementwurf} präsentiert Systementwurf und Architektur. Das Gesamtkonzept dieser sowie der Aufbau und Verantwortlichkeiten der Teilkomponenten werden dargelegt. Ebenfalls wird das Datenbankdesign und das allgemeine Sicherheits- bzw. Datenschutzkonzept beleuchtet.

Das darauffolgende Kapitel dokumentiert die Implementierung der einzelnen Services. Besonderes Augenmerk liegt hier auf der Integration der externen Komponenten wie Kameradaten oder der Verarbeitung der Bilddaten und des Datenflusses bei der Datenbeschaffung und -verarbeitung.

Das 6. Kapitel behandelt Infrastruktur, Deployment und Betrieb. Es wird auf die Container-Orchestrierung, CI/CD-Pipelines mit Deployment, Monitoring-Strategien sowie die Backup- und Restore-Strategie eingegangen. Auch wird kurz die MkDocs-basierte Entwicklerdokumentation erwähnt.

Kapitel~\ref{sec:qualitaetssicherung} widmet sich der Qualitätssicherung durch Tests und der Funktionsweise der Kennzeichenerkennung unter realen Bedingungen.

Das letzte Kapitel schließt mit einer Zusammenfassung und Evaluierung der Ergebnisse, der kritischen Reflexion des Projektverlaufs und einem Ausblick auf mögliche Erweiterungen.
